% ===== 6 工程部署与边界 =====
\section{现场部署、误差边界与改进优先级}
\label{sec:r-deploy}

\subsection{现场系统与硬件条件}

现场部署的首要目标不是追求单模块极限速度，而是保证\textbf{几何可追溯、视觉可质检、参数可冻结、结果可复现}。因此，硬件选择不以某个品牌或型号为核心，而以是否能够稳定恢复算法所依赖的物理条件为核心。固定俯视相机需要为 5--40~mm 主粒级提供足够像素覆盖；低畸变镜头和刚性支架保证一次测试内成像几何不漂移；深色哑光承载面与漫射照明降低硬阴影、高光和背景纹理造成的实例边界错误；与骨料测量面尽量共面的标尺或平面标记提供本次尺度锚点；CUDA GPU 则主要用于压缩 Cellpose-SAM 推理时间，使现场时间余量集中在标定、质检与失败恢复上。

表\ref{tab:r-hardware}给出适合本方案的工程配置及其与最终误差的关系。约 12~MP、固定焦距镜头等数值是当前推荐条件，而不是方法成立的唯一型号要求。若现场更换相机或镜头，应重新验证最小颗粒像素覆盖、透视与畸变，而不能仅因为“像素总数更高”就默认测量更准确。对静态自然堆积场景，普通固定曝光相机即可；若扩展到输送带，则全局快门和更短曝光优先级显著提高，因为运动模糊会直接改变轮廓几何。

\begin{table}[t]
\centering
\caption{推荐硬件条件及其与算法误差的直接关系。}
\label{tab:r-hardware}
\footnotesize
\renewcommand{\arraystretch}{1.12}
\begin{tabular}{>{\raggedright\arraybackslash}p{2.0cm}>{\raggedright\arraybackslash}p{4.0cm}>{\raggedright\arraybackslash}p{6.6cm}}
\toprule
模块 & 推荐条件 & 影响与现场验收重点 \\
\midrule
相机 & 约 12~MP 或更高；固定曝光；动态场景优先全局快门 & 保证 5~mm 附近仍有足够像素覆盖；检查无明显运动模糊、过曝或欠曝 \\
镜头/支架 & 固定焦距、低畸变、刚性俯视安装 & 保持一次测试内 mm/px 稳定；焦距、高度改变后必须重新标定 \\
承载面 & 深色、哑光、尽量平整 & 降低背景纹理和高光对边界的干扰；减少参考面与骨料面高度差 \\
照明 & 大面积漫射、亮度稳定 & 减少接触颗粒之间的硬阴影连接，降低 merge 风险 \\
尺度参考 & 标尺或已知尺寸平面标记，与测量面尽量共面 & 建立 pixel$\rightarrow$mm 锚点；多位置读数应基本一致 \\
计算设备 & 支持 CUDA 的 NVIDIA GPU & 缩短预训练模型推理并为现场重拍/复跑留出时间余量 \\
\bottomrule
\end{tabular}
\end{table}

软件链分为两个阶段。第一阶段由 \texttt{imagegrains} 完成实例分割、几何测量和毫米尺度写入，输出 mask、实例可视化及逐颗粒 CSV；第二阶段由 \texttt{aggregate\_screening} 读取逐颗粒表，计算筛分等效粒径、质量级配、形貌、异常与场景级报告。两阶段解耦后，如果上游实例已经通过质检，修改报告或重算下游统计不需要再次执行神经网络推理。正式比赛使用的模型版本、$\boldsymbol{\theta}$、$\gamma$ 和规则阈值在赛前冻结，现场只允许根据本次相机几何重新测量尺度 $r$。

\subsection{30 min 操作闭环与分层质检}

表\ref{tab:r-sop}给出面向 30~min 约束的操作预算。该表是工程上限而不是把历史 27~s/张算法时间简单外推为最终结论；正式赛前应完整演练并记录真实 wall-clock time。流程中特意保留重拍/复跑余量，因为尺度标错或整片 merge 一旦进入下游，会同时污染多个输出，现场尽早失败并恢复比事后解释异常数字更可靠。

\begin{table}[t]
\centering
\caption{评委版保留的 30~min 现场 SOP。}
\label{tab:r-sop}
\footnotesize
\renewcommand{\arraystretch}{1.12}
\begin{tabular}{>{\centering\arraybackslash}p{1.0cm}>{\raggedright\arraybackslash}p{3.2cm}>{\centering\arraybackslash}p{1.6cm}>{\raggedright\arraybackslash}p{7.2cm}}
\toprule
步骤 & 操作 & 预算 & 完成判据 \\
\midrule
1 & 设备就位、视场与照明 & 4~min & 目标完整入镜、相机固定、无明显硬阴影/模糊 \\
2 & 尺度标定 & 3~min & 获得本场景 mm/px；多位置参考无明显冲突 \\
3 & 采集与快速质检 & 3~min & 图像清晰、参考物共面且可读取 \\
4 & 分割与几何测量 & 5~min & 生成实例 mask/CSV；无大面积 merge 或数量级异常 \\
5 & 筛分、形貌、异常分析 & 2~min & 生成场景 summary、级配曲线和可视化 \\
6 & 结果质检与必要复跑 & 5~min & 尺度、轴长、实例和级配量级一致；必要时重拍/复跑固定配置 \\
7 & 导出与核对 & 3~min & 参数、逐颗粒表、场景报告和图像同时留存 \\
-- & 故障余量 & 5~min & 应对设备连接、模型加载或一次重拍 \\
\bottomrule
\end{tabular}
\end{table}

现场质检被设计为三个逐层门槛。\textbf{采集门槛}在神经网络推理前检查视场、曝光、清晰度、参考物和严重遮挡；这一层失败时直接重拍，因为输入证据本身已经不足。\textbf{实例门槛}在推理后检查 composite/overlay，重点寻找大面积 merge、异常碎裂、边缘截断和明显漏检；这一层失败时首先判断原因是否来自图像质量，若是则回到采集，否则只能按预先冻结的备用分割配置复跑，不能根据最终 $D_{50}$ 是否“好看”来临时调模型。\textbf{物理/报告门槛}检查参考长度反算、轴长量级、异常比例、筛级总和和输出文件完整性；这一层失败时停止提交，先定位尺度或数据链错误。

表\ref{tab:r-recovery}把几类最可能发生的现场异常转换为明确恢复动作。这里的核心原则是\textbf{根据中间证据决定是否重拍/复跑，而不是根据期望答案反向调参}。如果只看到最终级配曲线而没有实例和尺度证据，就无法区分“样品本来偏粗”和“模型整片 merge”；因此中间可视化不是展示附件，而是现场质量控制的一部分。

\begin{table}[t]
\centering
\caption{现场失败信号、原因判别与恢复动作。}
\label{tab:r-recovery}
\footnotesize
\renewcommand{\arraystretch}{1.11}
\begin{tabular}{>{\raggedright\arraybackslash}p{3.0cm}>{\raggedright\arraybackslash}p{3.7cm}>{\raggedright\arraybackslash}p{4.1cm}>{\raggedright\arraybackslash}p{2.3cm}}
\toprule
失败信号 & 优先检查原因 & 恢复动作 & 是否允许改模型参数 \\
\midrule
参考物模糊/尺度读数冲突 & 焦点、共面性、相机倾斜、镜头畸变 & 重拍或调整姿态；重新计算 $r$ & 否 \\
大面积整片 merge & 硬阴影、曝光、严重接触或域失配 & 优先改善采集并重拍；仍失败再用赛前冻结备用配置 & 仅冻结方案 \\
大量细碎 split & 分辨率、纹理、高光、后处理下限 & 检查图像质量；按冻结最小实例规则复跑 & 不临场搜索 \\
轴长方向/量级明显异常 & 错 CSV、旧 resolution、截断实例 & 核对 batch id 与尺度；必要时重跑测量 & 否 \\
$D_p$ 与场景观感突变 & merge、极端异常、质量权重输入错误 & 向上检查实例和逐颗粒表，不因结果预期改 $\gamma$ & 否 \\
输出文件缺失/不一致 & 路径、上一批文件混用、流程中断 & 从最近通过质检的中间产物重算并核对批次 id & 否 \\
\bottomrule
\end{tabular}
\end{table}

这一门控式流程还定义了“何时应承认本次单目测量无效”。如果自然堆积状态导致大部分颗粒严重遮挡，或参考物无法与测量面建立可信尺度关系，继续运行模型只会产生形式完整但缺乏物理依据的数字。此时正确的工程行为不是无限后处理，而是记录失败原因并重新采集；若比赛场景本身不允许改善这些条件，则应把它归入方法边界，在后续方案中用多视角或深度传感器解决。

\subsection{结果输出与可追溯证据链}

比赛现场最终不应只导出一张“结果图”。为了使评委能够从场景级指标反查到具体颗粒和运行参数，系统保留四层输出：原始输入与尺度、实例 mask/可视化、逐颗粒结构化表、场景级汇总。表\ref{tab:r-output-chain}给出最小提交/留档内容。每个场景使用唯一 batch/scene id 关联这些文件，使 $D_{50}$、某个筛级比例或异常数量都能够回溯到对应的 $d_i$、mask 和原始图像。

\begin{table}[t]
\centering
\caption{场景级结果包及其审计作用。}
\label{tab:r-output-chain}
\footnotesize
\renewcommand{\arraystretch}{1.12}
\begin{tabular}{>{\raggedright\arraybackslash}p{2.7cm}>{\raggedright\arraybackslash}p{5.0cm}>{\raggedright\arraybackslash}p{5.4cm}}
\toprule
层级 & 保存内容 & 作用 \\
\midrule
输入层 & 原始图像、场景 id、尺度参考与 $r$ & 重新检查拍摄条件和毫米尺度来源 \\
实例层 & integer mask、detection/composite、axes overlay & 复核 merge/split、漏检、轴长方向与实例定位 \\
颗粒层 & $a,b,A,P,d_{\mathrm{eq}},d,w$、形貌与异常标签 & 复算任一筛级或 $D_p$，定位极端质量权重来源 \\
场景层 & $D_{10}/D_{50}/D_{90}$、逐筛占比、形貌/异常汇总、级配曲线 & 形成比赛主结果与答辩展示 \\
配置层 & 模型版本、$\boldsymbol{\theta}$、$\gamma$、规则版本、运行时间 & 保证同一结果能够在冻结配置下复现 \\
\bottomrule
\end{tabular}
\end{table}

正式机械筛分验证阶段还应为同一 batch id 增加原始称量记录，包括总质量、各筛上质量、累计通过率和试验日期。这样“原图--实例--预测级配--机械真值”才能形成完整闭环，也便于在出现某一筛级大误差时判断问题来自视觉实例、筛分粒径映射还是质量模型，而不是只得到一个无法诊断的总误差。

\subsection{误差传播与方法边界}

视觉筛分的最终误差可理解为尺度、实例、筛分粒径映射和质量模型误差共同传播的结果。它们的重要性并不相同：尺度偏差会系统性移动全部粒径；merge 在 $d^{\gamma}$ 下会放大粗粒质量；遮挡和第三维缺失则属于单目观测本身的不可辨识信息。表\ref{tab:r-limits}按对比赛主要筛分指标的影响优先级总结当前边界和改进路线。

\begin{table}[t]
\centering
\caption{主要误差源、对最终级配的传播路径与改进优先级。}
\label{tab:r-limits}
\footnotesize
\renewcommand{\arraystretch}{1.13}
\begin{tabular}{>{\raggedright\arraybackslash}p{2.7cm}>{\raggedright\arraybackslash}p{4.1cm}>{\raggedright\arraybackslash}p{4.3cm}>{\centering\arraybackslash}p{1.5cm}}
\toprule
误差源 & 对结果的主要影响 & 当前缓解/下一步 & 优先级 \\
\midrule
缺少配对机械筛分真值 & 无法判断 $\boldsymbol{\theta},\gamma$ 是否逼近标准筛分 & 同批图像--机械筛分；校准/测试物理隔离 & 最高 \\
尺度与透视 & 全场 $D_p$ 与筛级边界系统偏移 & 共面参考、多位置检查；畸变/单应校正 & 高 \\
merge/split 与域迁移 & merge 偏粗且被质量权重放大；split 偏细 & 人工 mask 困难集；标准成像后再决定是否微调 & 高 \\
遮挡与不可见颗粒 & 局部轮廓低估或完全无观测证据 & 降低严重二层堆叠；必要时多视角/深度 & 高 \\
第三维厚度缺失 & 二维短轴与真实筛孔通过行为存在残差 & 低维物理校准；残差稳定时再引入深度/多视角 & 中高 \\
质量幂律近似 & 不同形貌/尺寸的真实单颗粒质量不完全服从同一 $\gamma$ & 多批次联合校准；后续加入厚度/形貌项 & 中 \\
泥团语义 & 几何阈值可能漏检或误报材质异常 & 目前仅输出疑似候选；后续颜色/纹理分类器 & 中低 \\
\bottomrule
\end{tabular}
\end{table}

其中最高优先级并不是继续增加无真值截图或更换更大的视觉模型，而是获取能直接约束比赛评分对象的\textbf{同批机械筛分 ground truth}。如果下游物理语义尚未验证，即使上游 mask 看起来非常精细，也无法证明最终标准筛分误差足够小。第二优先级是稳定成像和实例边界，因为它们影响全部后续测量；只有在标准化采集后仍出现系统性分割失败，才值得用本赛题少量人工 mask 微调模型。

单目方案还有两项不能通过“继续调参数”消除的物理边界。第一，完全被遮挡的颗粒没有成像证据，任何恢复都依赖统计猜测；第二，二维投影缺少厚度，因此不能从理论上唯一恢复三维针片状或每个颗粒的真实质量。如果赛题现场的堆叠和三维形貌误差最终成为主导项，更合理的升级是多视角、深度相机或受控铺料，而不是无限增加二维后处理复杂度。

综合而言，当前路线适合比赛阶段的原因在于它把系统复杂度集中在真正有证据收益的地方：上游利用成熟预训练实例感知，下游用少量机械筛分数据校准赛题特有的物理/统计关系，现场保持单目、快速标定和固定参数推理。该设计也为后续升级留出清晰接口，而不要求推翻整条处理链。
